% =============================================================================
% Coordinate system pic (for pendulum hinge)
% =============================================================================
\tikzset{
  coordsys/.pic={
      \fill[red] (0,0) circle(0.05);
      \draw[->,thick, red] (0,0) -- (1,0) node[right] {$x$};
      \draw[->,thick, red] (0,0) -- (0,1) node[above] {$y$};
    }
}

% =============================================================================
% Pendulum with wall contact
% =============================================================================
\newcommand{\PendulumComponent}{
\def\L{4}       % rod length
\def\W{0.3}     % rod width
\def\R{0.5}     % pendulum head radius
\def\ang{40}    % angle from negative y-axis (hanging down)
\def\F{1.6}     % force arrow magnitude
\def\wallW{0.5} % wall thickness
\def\wallH{2}   % wall height

\coordinate (O) at (0,0);

\draw[thick, fill=lightgray!20] (-1.5, -\L-\wallH/2-0.5) rectangle (7.5, 2.5);
\draw node at (3, 1.5) [above, font=\large\bfseries] {Pendulum with Wall Contact};

% Pendulum rod (rectangle, rotated about hinge)
\begin{scope}[rotate=\ang]
  \draw[thick, fill=brown!30] (-\W/2,0) rectangle (\W/2,-\L);
  % Pendulum head (circle at end of rod)
  \draw[thick, fill=gray!40] (0,-\L) circle(\R);
  \node[right=\R+12pt] at (0,-\L) {$m,\,I$};
  % Rod midline (dash-dot) — extended beyond head for force reference
  \draw[dash dot, thin] (0,0.3) -- (0,-\L-\R-1.2);
  % Pendulum head center lines (dash-dot, cross at center of circle)
  \draw[dash dot, thin] (-\R-0.3,-\L) -- (\R+0.3,-\L);
  % Origin center line parallel to pendulum head center line (dash-dot)
  \draw[dash dot, thin] (-\R-0.3,0) -- (\R+0.3,0);
\end{scope}

% === Ghost pendulum at theta = 0 (contact position) ===
\begin{scope}[opacity=0.35]
  \draw[thick, fill=brown!15] (-\W/2,0) rectangle (\W/2,-\L);
  \draw[thick, fill=gray!20] (0,-\L) circle(\R);
\end{scope}

% Vertical reference line (theta = 0)
\draw[dash dot, thin] (0,0) -- (0,-\L-\R-0.3);

% Pendulum trajectory (arc of the head center)
\draw[dotted, thick] (-90:\L) arc (-90:-90+\ang+0:\L);

% Wall (impact surface at theta = 0, right edge at x = -R)
\draw[supportSurface] (-\R-\wallW, -\L-\wallH/2) rectangle (-\R, -\L+\wallH/2);
\draw[thick] (-\R, -\L-\wallH/2) -- (-\R, -\L+\wallH/2);

% === Forces at pendulum head ===
\coordinate (M) at (\ang-90:\L);

% Gravity force (mg, straight down from head center)
\coordinate (FG) at ($(M)+(-90:{\F+\R})$);
\draw[forceVector] (M) -- (FG) node[right=1pt] {$mg$};

% Radial component of gravity (along rod, away from hinge)
\coordinate (FGrad) at ($(M)+(-90+\ang:{\F*cos(\ang)+\R})$);
\draw[forceVector, forceComponentColor] (M) -- (FGrad) node[right=1pt] {$mg\cos\theta$};

% Tangential component of gravity (perpendicular to rod)
\coordinate (FGtan) at ($(M)+(180+\ang:{\F*sin(\ang)+\R})$);
\draw[forceVector, forceComponentColor] (M) -- (FGtan) node[above=15pt, left=-18pt] {$mg\sin\theta$};

% Projection dashed lines (parallelogram construction)
\draw[dashed, forceColor!40] (FGtan) -- (FG);
\draw[dashed, forceColor!40] (FGrad) -- (FG);

% Theta angle at pendulum head (between gravity and radial direction)
\draw pic[-{Latex[length=3,width=2]},thin,"$\theta$",angleColor,draw=angleColor,
angle radius=14, angle eccentricity=1.5] {angle=FG--M--FGrad};

% Gravity direction indicator (general, on the left side)
\draw[->, very thick, black!60] (-1,-0.5) -- (-1,-1.5) node[right=2pt] {$g$};

% Angle theta at hinge (between vertical and rod midline)
\coordinate (B) at (0,-1.5);
\draw pic[->,"$\theta$", draw, angle radius=60, angle eccentricity=1.15] {angle=B--O--M};

% Length L annotation (offset so it clears the hinge circle)
\draw[<->,thick] (\ang-90+90:0.7) ++ (\ang-90:0) -- ++ (\ang-90:\L)
node[midway, above, sloped] {$L$};

% Revolute joint: classical triangle + small circle
\draw[thick, fill=white] (-0.45,0.6) -- (0.45,0.6) -- (0,0) -- cycle;
\draw[thick, fill=white] (O) circle(\W/2);

% Drive torque at hinge
\draw[<-, very thick, blue!70!black] (160:0.5) arc (160:70:0.5)
node[left=15pt] {$\tau$};

% Coordinate system at hinge (topmost layer)
\pic at (0,0) {coordsys};

\draw[xshift=5.5cm, yshift=-0.75cm]
node [text width=3cm, information text]
{Dynamics:
  \begin{align*}
    \dot{\theta} & = \omega                           \\
    \dot{\omega} & = \frac{\tau - mgL\sin(\theta)}{I}
  \end{align*}
  Rigid contact:
  \begin{align*}
    \omega^+ &= -\omega^-
  \end{align*}
};

% --- Interface anchors (for wiring in assembly) ---
\coordinate (pend-tau)   at (0, 0);     % torque input
\coordinate (pend-theta) at (0, 0);     % angle output (same shaft)
}